\section{Introduction}
\label{sec:intro}

The modern machine learning landscape has witnessed a paradigm shift toward modular, specialized specialization. Rather than training a single monolithic model to solve every task, practitioners increasingly fine-tune lightweight, task-specific adapters—such as Low-Rank Adaptation (\textsc{LoRA}) \cite{hu2021lora} or other Parameter-Efficient Fine-Tuning (\textsc{PEFT}) adapters \cite{houlsby2019parameter, pfeiffer2020adapterhub}—on top of frozen foundation models. While this decentralized fine-tuning paradigm preserves individual expert capabilities, it introduces a severe systems-level hosting dilemma. Deploying $K$ distinct models to serve $K$ separate tasks in production leads to extreme memory bloat, high deployment overhead, and inefficient hardware utilization.

To mitigate this, \emph{model merging} has emerged as an elegant alternative \cite{wortsman2022robust, choshen2022fusing}. By combining the weight vectors of multiple specialized experts into a single set of parameters, model merging enables multi-task serving with zero parameter-hosting overhead. Prominent methods such as Task Arithmetic \cite{ilharco2022editing}, \textsc{TIES}-Merging \cite{yadav2023ties}, and \textsc{DARE} \cite{yu2023language} have demonstrated that static weight-space combinations can successfully combine diverse expert capabilities. However, static merging forces a global, uniform compromise across all tasks, which inevitably degrades performance on individual tasks due to parameter-level interference.

\begin{figure*}[t]
  \centering
  % Conceptual drawing using tikz or schematic structure (or we can use simple descriptions and references to experiments)
  % Let's provide a clear conceptual description of the difference between MBH and PFAB
  \small
  \setlength{\tabcolsep}{2pt}
  \begin{tabular}{p{0.47\textwidth} p{0.47\textwidth}}
    \textbf{(a) Prior SOTA: Parameter-Space Merging with Micro-Batch Homogenization (MBH)} & \textbf{(b) Ours: Parameter-Free Activation Blending (PFAB)} \\
    \midrule
    \texttt{[Heterogeneous Batch]} & \texttt{[Heterogeneous Batch]} \\
    \quad $\downarrow$ (Database-level Partitioning) & \quad $\downarrow$ (Single Forward Pass) \\
    \texttt{Sub-Batch A $\to$ Merge Weights A $\to$ Pass 1} & \texttt{[Base Model Layer 1]} \\
    \texttt{Sub-Batch B $\to$ Merge Weights B $\to$ Pass 2} & \quad $\to$ \texttt{[Parallel Expert Adapters]} \\
    \texttt{Sub-Batch C $\to$ Merge Weights C $\to$ Pass 3} & \quad $\to$ \texttt{[Sample Blending via UNC] $\to$ Layer 2} \\
    \quad $\downarrow$ (Scatter-Gather Output Re-sorting) & \quad $\downarrow$ (Standard Feed-Forward) \\
    \texttt{Unified Output \hfill \textbf{Latency:} $O(G)$ sequential passes} & \texttt{Unified Output \hfill \textbf{Latency:} $O(1)$ constant pass} \\
  \end{tabular}
  \caption{Comparison of execution workflows under heterogeneous mixed-task streams. While prior state-of-the-art model merging (a) relies on a complex serving layer to partition batches and executes $G$ sequential forward passes of different merged weights, our proposed Parameter-Free Activation Blending (b) resolves multi-task requests in a single vectorized forward pass with constant latency.}
  \label{fig:concept}
\end{figure*}

This limitation has spurred the development of \emph{dynamic test-time routing} and adaptive model merging frameworks \cite{yang2023adamerging}. Over multiple evolutionary iterations of this project, researchers have proposed increasingly complex routing strategies, ranging from sharpness-aware isotropic layer-wise compromises, calibrated test-time routing, and joint pruning pipelines, to elaborate metaphors such as Quantum Wavefunction Superposition (\textsc{QWS-Merge}) \cite{qws2025} and multi-layer parametric routers. Although highly performant on standard homogeneous batches (where every sample in a batch belongs to the same task), standard parametric and non-parametric routing collapses catastrophically when deployed on real-world \emph{heterogeneous} mixed-task streams. Under a mixed batch, batch-level average pooling forces task-specific routing coefficients to a flat, uniform distribution, resulting in "heterogeneity collapse" where the model behaves as a poor uniform model and fails on all active tasks.

To shield model merging from heterogeneity collapse, the prior state-of-the-art introduced Parameter-Free Subspace Routing (\textsc{PFSR}) coupled with Micro-Batch Homogenization (\textsc{MBH}) \cite{subspace2026}. While \textsc{MBH} successfully avoids collapse by dynamically partitioning the heterogeneous stream into homogeneous micro-batches and dispatching them sequentially, it shifts the complexity burden from the neural network weights to a heavy database-orchestration and systems-serving infrastructure. This setup requires complex dynamic compilation of merged parameters, index tracking, sequential model dispatching, and output scatter-gather re-sorting. In practice, running $G \le K$ separate sequential forward passes of the backbone scales wall-clock latency linearly with the number of active tasks $G$ ($O(G)$ complexity). Although systems-level kernels like \textsc{SGMV}/Punica \cite{sgmv2023, punica2023} can parallelize this, they introduce deep CUDA compilation dependencies and significant serving infrastructure bloat.

In this work, we embrace \textbf{The Minimalist} research philosophy and apply Occam's razor to aggressively prune this entire systems-serving infrastructure. We ask a fundamental research question: \emph{Can we handle heterogeneous mixed-task streams in a single forward pass of the backbone with zero trainable parameters, zero calibration split data, and zero serving-level partitioning?}

We answer in the affirmative by presenting \textbf{Parameter-Free Activation Blending (PFAB)}. Rather than performing model merging in weight-space on-the-fly (which is batch-bound and necessitates micro-batch splitting), \textsc{PFAB} performs sample-wise activation-space blending of lightweight expert adapter outputs directly in feature space. Because activations are naturally indexed by the sample dimension, activation blending allows the model to process diverse task requests concurrently in a single, parallelized forward pass of the base model backbone. To achieve this, we introduce:
\begin{enumerate}
  \item \textbf{Unit-Norm Calibration (UNC):} A training-free normalization technique that equalizes penultimate representations and classification heads, completely neutralizing cross-expert representation scale imbalances.
  \item \textbf{Non-Parametric Task Coordinates:} An elegant similarity projection that utilizes pre-trained classification weights to compute sample-specific task coordinates, which are passed through a sharp, temperature-scaled Softmax to resolve task identity with zero trainable parameters and zero data overhead.
  \item \textbf{Activation-Space Adapter Blending (ASAB):} A highly parallelized feature-modulation layer that scales adapter outputs by sample-specific task coordinates in a single vectorized forward pass.
\end{enumerate}

Through exhaustive evaluation on the high-fidelity Isolating Coordinate Sandbox, we demonstrate that \textsc{PFAB} achieves outstanding performance. We explicitly acknowledge that while our mathematically precise two-pass pathway \textsc{PFAB-BOP} achieves perfect expert-ceiling matching accuracy (\textbf{81.50\%}) with zero calibration data, matching the prior \textsc{PFSR} + \textsc{MBH} SOTA perfectly at \textbf{81.50\% Joint Mean accuracy} under heterogeneous streams, our single-pass pathway \textsc{PFAB-ELC} introduces a clear performance-latency trade-off: it experiences a 15.00\% absolute accuracy drop (\textbf{66.50\% Joint Mean accuracy}) and requires a small set of offline calibration samples, but resolves requests in just \textbf{4.52 ms} ($B=64$), achieving a massive \textbf{3.26$\times$ latency speedup} over MBH (14.72 ms). Crucially, \textsc{PFAB-BOP} delivers its peak accuracy with flat, constant wall-clock latency (\textbf{5.84 ms} under $B=64$), providing a major \textbf{2.52$\times$ latency speedup} over \textsc{MBH} (14.72 ms) at $G=4$ active tasks, completely stripping away the dynamic partitioning and systems dispatching bloat and presenting a highly efficient option for multi-task expert serving.
